
\subsubsection{\textcolor{red}{周五}耶穌被馬利亞所膏}

\begin{table}[H]
  \centering
    \begin{tabular}{p{14em}|p{7.93em}|l|l|l|p{4.855em}}
    \toprule
    \textbf{事 件} & \textbf{地 點} & \multicolumn{1}{p{5.285em}|}{\textbf{馬太福音}} & \multicolumn{1}{p{5.57em}|}{\textbf{馬可福音}} & \multicolumn{1}{p{5.285em}|}{\textbf{路加福音}} & \textbf{約翰福音} \\  \midrule

\textcolor{red}{周五}耶穌被馬利亞所膏 & 伯大尼西門家& \multicolumn{1}{p{5.285em}|}{26:6-13} & \multicolumn{1}{p{5.57em}|}{14:3-9} &       & 12:1-8 \\ \midrule

祭司長謀算要殺主和拉撒路 & 伯大尼？   &       &       &       & 12:9-11 \\\midrule
\end{tabular}
\end{table}

\begin{itemize}
\item 馬太 26:6-13\\
耶穌在伯大尼長大痲瘋的西門家裡，有一個女人拿著一玉瓶極貴的香膏來，趁耶穌坐席的時候，澆在他的頭上。門徒看見就很不喜悅，說：何用這樣的枉費呢！這香膏可以賣許多錢，賙濟窮人。耶穌看出他們的意思，就說：為什麼難為這女人呢？他在我身上做的是一件美事。因為常有窮人和你們同在；只是你們不常有我。他將這香膏澆在我身上是為我安葬做的。我實在告訴你們，普天之下，無論在什麼地方傳這福音，也要述說這女人所行的，作個紀念。
\item 馬可 14:3-9\\
只是說：當節的日子不可，恐怕百姓生亂。
耶穌在伯大尼長大痲瘋的西門家裡坐席的時候，有一個女人拿著一玉瓶至貴的真哪噠香膏來，打破玉瓶，把膏澆在耶穌的頭上。
有幾個人心中很不喜悅，說：何用這樣枉費香膏呢？
這香膏可以賣三十多兩銀子賙濟窮人。他們就向那女人生氣。
耶穌說：由他吧！為什麼難為他呢？他在我身上做的是一件美事。
因為常有窮人和你們同在，要向他們行善隨時都可以；只是你們不常有我。
他所做的，是盡他所能的；他是為我安葬的事把香膏預先澆在我身上。
我實在告訴你們，普天之下，無論在什麼地方傳這福音，也要述說這女人所做的，以為記念。
\item 約翰 12:1-8\\
逾越節前六日，耶穌來到伯大尼，就是他叫拉撒路從死裡復活之處。
有人在那裡給耶穌預備筵席；馬大伺候，拉撒路也在那同耶穌坐席的人中。
馬利亞就拿著一斤極貴的真哪噠香膏，抹耶穌的腳，又用自己頭髮去擦，屋裡就滿了膏的香氣。有一個門徒，就是那將要賣耶穌的加略人猶大，
說：這香膏為什麼不賣三十兩銀子賙濟窮人呢？
他說這話，並不是掛念窮人，乃因他是個賊，又帶著錢囊，常取其中所存的。
耶穌說：由他吧！他是為我安葬之日存留的。
因為常有窮人和你們同在，只是你們不常有我。
\item 祭司長等要除滅耶穌-約翰 12:9-11\\
有許多猶太人知道耶穌在那裡，就來了，不但是為耶穌的緣故，也是要看他從死裡所復活的拉撒路。但祭司長商議連拉撒路也要殺了；因有好些猶太人為拉撒路的緣故，回去信了耶穌。
\end{itemize}

\subsubsection{\textcolor{red}{周六}騎驢駒進耶京合城驚動}

\begin{table}[H] \centering
 \begin{tabular}{p{14em}|p{7.93em}|l|l|l|p{4.855em}} \toprule
\textbf{事 件} & \textbf{地 點} & \multicolumn{1}{p{5.285em}|}{\textbf{馬太福音}} & \multicolumn{1}{p{5.57em}|}{\textbf{馬可福音}} & \multicolumn{1}{p{5.285em}|}{\textbf{路加福音}} & \textbf{約翰福音} \\  \midrule
打發兩個門徒牽驢駒 &\textcolor{red}{周六}橄欖山&21:1-5 &  11:1-3   &19:28-31&  \\\midrule
門徒如耶穌所說牽驢駒 &伯法其和伯大尼& 21:6&  11:4-6     &19:32-34&  \\\midrule

門徒搭衣在驢駒上鋪路上&伯法其和伯大尼& 21:7-8&  11:7-8     &19:35-36& \\\midrule

眾門徒眾人歡樂大聲讚美神&\textcolor{red}{周六晚}下橄欖山& 21:9&11:9-10 &19:37-38&12:12-16\\\midrule
見拉撒路復活的眾人作見證&下橄欖山& &&&12:17-18\\\midrule
法利賽人彼此叹世人都隨主&下橄欖山&       &       && 12:19\\\midrule
法利賽人劝耶穌責備門徒&下橄欖山&       &       &19:39-40&  \\\midrule

為耶路撒冷哀哭预其将被毁 &快到耶路撒冷& 23:37-39 &       &19:41-44&  \\\midrule
耶穌既進耶京合城驚動&耶路撒冷& 21:10-11&   &&  \\\midrule
腓力安得烈同幾個希臘人見主&下橄欖山&       &       &&12:20-22\\\midrule

入聖殿周圍看了各樣物件&耶路撒冷& & 11:11a&&  \\\midrule
天色已晚和門徒出城&伯大尼& &11:11b &&  \\\midrule
\end{tabular}%
\end{table}%

\begin{itemize}
\item 耶穌騎驢進京
\begin{itemize}
\item 馬太 21:1-11\\
耶穌和門徒將近耶路撒冷，到了伯法其，在橄欖山那裡。耶穌就打發兩個門徒，對他們說：你們往對面村子裡去，必看見一匹驢拴在那裡，還有驢駒同在一處；你們解開，牽到我這裡來。若有人對你們說什麼，你們就說：主要用他。那人必立時讓你們牽來。這事成就是要應驗先知的話，說：要對錫安的居民（原文是女子）說：看哪，你的王來到你這裡，是溫柔的，又騎著驢，就是騎著驢駒子。門徒就照耶穌所吩咐的去行，牽了驢和驢駒來，把自己的衣服搭在上面，耶穌就騎上。眾人多半把衣服鋪在路上；還有人砍下樹枝來鋪在路上。前行後隨的眾人喊著說：和散那（原有求救的意思，在此是稱頌的話）歸於大衛的子孫！奉主名來的是應當稱頌的！高高在上和散那！耶穌既進了耶路撒冷，合城都驚動了，說：這是誰？
眾人說：這是加利利拿撒勒的先知耶穌。
\item 馬可 11:1-10\\
耶穌和門徒將近耶路撒冷，到了伯法其和伯大尼，在橄欖山那裡；耶穌就打發兩個門徒，
對他們說：你們往對面村子裡去，一進去的時候，必看見一匹驢駒拴在那裡，是從來沒有人騎過的，可以解開，牽來。若有人對你們說：為什麼做這事？你們就說：主要用他。那人必立時讓你們牽來。他們去了，便看見一匹驢駒拴在門外街道上，就把他解開。在那裡站著的人，有幾個說：你們解驢駒做什麼？門徒照著耶穌所說的回答，那些人就任憑他們牽去了。
他們把驢駒牽到耶穌那裡，把自己的衣服搭在上面，耶穌就騎上。有許多人把衣服鋪在路上，也有人把田間的樹枝砍下來，鋪在路上。前行後隨的人都喊著說：和散那（和散那：原有求救的意思，在此乃是稱頌的話）！奉主名來的是應當稱頌的！那將要來的我祖大衛之國是應當稱頌的！高高在上和散那！
\item 路加 19:28-40\\
将近伯法其和伯大尼，在一座山名叫橄榄山那里。就打发两个门徒，说， 你们往对面村子里去。进去的时候，必看见一匹驴驹拴在那里，是从来没有人骑过的。可以解开牵来。 若有人问为什么解它，你们就说，主要用它。 打发的人去了，所遇见的，正如耶稣所说的。 他们解驴驹的时候，主人问他们说，解驴驹作什么？ 他们说，主要用它。 他们牵到耶稣那里，把自己的衣服搭在上面，扶着耶稣骑上。 走的时候，众人把衣服铺在路上。 将近耶路撒冷，正下橄榄山的时候，众门徒因所见过的一切异能，都欢乐起来，大声赞美神， 说，奉主名来的王，是应当称颂的。在天上有和平，在至高之处有荣光。 众人中有几个法利赛人对耶稣说，夫子，责备你的门徒吧。 耶稣说，我告诉你们，若是他们闭口不说，这些石头必要呼叫起来。
\item 約翰 12:12-19\\
第二天，有許多上來過節的人聽見耶穌將到耶路撒冷，
就拿著棕樹枝出去迎接他，喊著說：和散那！奉主名來的以色列王是應當稱頌的！
耶穌得了一個驢駒，就騎上，如經上所記的說：
錫安的民（原文作女子）哪，不要懼怕！你的王騎著驢駒來了。
這些事門徒起先不明白，等到耶穌得了榮耀以後才想起這話是指著他寫的，並且眾人果然向他這樣行了。
當耶穌呼喚拉撒路，叫他從死復活出墳墓的時候，同耶穌在那裡的眾人就作見證。
眾人因聽見耶穌行了這神蹟，就去迎接他。
法利賽人彼此說：看哪，你們是徒勞無益，世人都隨從他去了。
\item 综合\\
耶穌和門徒將近耶路撒冷，將近伯法其和伯大尼，到了伯法其，在橄欖山那裡， 耶穌就打發兩個門徒，對他們說：「你們往對面村子裡去，必看見一匹驢拴在那裡，還有驢駒同在一處, 是從來沒有人騎過的。你們解開，牽到我這裡來。若有人對你們說什麼，你們就說：『主要用牠。』那人必立時讓你們牽來。」 這事成就，是要應驗先知的話說： 「要對錫安的居民 說：『看哪，你的王來到你這裡，是溫柔的，又騎著驢，就是騎著驢駒子。』」門徒就照耶穌所吩咐的去行，所遇見的正如耶穌所說的, 便看見一匹驢駒拴在門外街道上，就把牠解開。他們解驢駒的時候，在那裡站著的人有幾個, 主人問他們說：「解驢駒做什麼？」他們說：「主要用牠。」那些人就任憑他們牽去了。 牽了驢和驢駒來，把自己的衣服搭在上面，扶著耶穌騎上, 耶穌就騎上。  眾人多半把衣服鋪在路上，還有人把田間的樹枝砍下樹枝來鋪在路上。將近耶路撒冷，正下橄欖山的時候，眾門徒因所見過的一切異能都歡樂起來，大聲讚美神， 說：「奉主名來的王是應當稱頌的！在天上有和平，在至高之處有榮光！」 前行後隨的眾人喊著說：「和散那歸於大衛的子孫！奉主名來的是應當稱頌的！那將要來的我祖大衛之國是應當稱頌的！高高在上和散那！」眾人中有幾個法利賽人對耶穌說：「夫子，責備你的門徒吧！」 耶穌說：「我告訴你們，若是他們閉口不說，這些石頭必要呼叫起來！」
\end{itemize}
\item 為耶路撒冷哀哭
\begin{itemize}
\item 馬太 23:37-39\\
耶路撒冷啊，耶路撒冷啊，你常殺害先知，又用石頭打死那奉差遣到你這裡來的人。我多次願意聚集你的兒女，好像母雞把小雞聚集在翅膀底下，只是你們不願意。
看哪，你們的家成為荒場留給你們。
我告訴你們，從今以後，你們不得再見我，直等到你們說：奉主名來的是應當稱頌的。
\item 路加 19:41-44\\
耶稣快到耶路撒冷看见城，就为他哀哭， 说，巴不得你在这日子，知道关系你平安的事。无奈这事现在是隐藏的，叫你的眼看不出来。 因为日子将到，你的仇敌必筑起土垒，周围环绕你，四面困住你， 并要扫灭你，和你里头的儿女，连一块石头也不留在石头上。因你不知道眷顾你的时候。 
\item 综合\\
耶穌快到耶路撒冷，看見城，就為它哀哭說：「巴不得你在這日子知道關係你平安的事，無奈這事現在是隱藏的，叫你的眼看不出來。 因為日子將到，你的仇敵必築起土壘，周圍環繞你，四面困住你，並要掃滅你和你裡頭的兒女，連一塊石頭也不留在石頭上，因你不知道眷顧你的時候。」
\end{itemize}
\item 查看聖殿-馬可11:11a\\
耶穌進了耶路撒冷，入了聖殿，周圍看了各樣物件。天色已晚，就和十二個門徒出城，往伯大尼去了。

\item 夜宿伯大尼
\begin{itemize}
\item 馬可11:11b - 往伯大尼去了。
\item 耶穌既進了耶路撒冷，合城都驚動了，說：「這是誰？」 眾人說：「這是加利利拿撒勒的先知耶穌。」耶穌入了聖殿，周圍看了各樣物件。天色已晚，就和十二個門徒出城，往伯大尼去了。
\end{itemize}
\end{itemize}

\begin{table}[htbp] \centering  
\begin{tabular}{l|l|l} \toprule
\multicolumn{3}{l}{1， 耶穌騎驢進京} \\\midrule
\multicolumn{1}{c|}{\multirow{3}{*}{在橄欖山那裡}} & 門徒    & \multicolumn{1}{|p{29.9em}}{把自己的衣服搭在上面，扶著耶穌騎上, 耶穌就騎上}  \\
\cmidrule{2-3}          & \multicolumn{1}{l|}{ 眾人} & \multicolumn{1}{|p{29.9em}}{眾人多半把衣服鋪在路上，還有人把田間的樹枝砍下樹枝來鋪在路上}  \\    \midrule
 & 眾門徒   & \multicolumn{1}{|p{29.9em}}{眾門徒都歡樂起來，大聲讚美神說：「奉主名來的王是應當稱頌的！在天上有和平，在至高之處有榮光！」}   \\
\cmidrule{2-3}          & \multicolumn{1}{l}{跟隨眾人} &\multicolumn{1}{|p{32.9em}}{ 前行後隨的眾人喊著說：「和散那歸於大衛的子孫！奉主名來的是應當稱頌的！那將要來的我祖大衛之國是應當稱頌的！高高在上和散那！」} \\    
\cmidrule{2-3}    \multicolumn{1}{p{7em}|}{下橄欖山的時候} & 法利賽人  &  \multicolumn{1}{|p{29.9em}}{眾人中有幾個法利賽人對耶穌說：「夫子，責備你的門徒吧！」} \\
\cmidrule{2-3}          & 耶穌    & \multicolumn{1}{|p{29.9em}}{「我告訴你們，若是他們閉口不說，這些石頭必要呼叫起來！」} \\   \midrule  
\multicolumn{3}{l}{2，耶穌為聖城哀哭} \\\midrule
\multicolumn{3}{l}{3，耶穌在聖城行程並夜宿伯大尼} \\\bottomrule
 \end{tabular}%
 \caption{週六}\label{tab:addlabel}%
\end{table}%




\subsubsection{周日}
\begin{table}[H]\centering
\begin{tabular}{p{14em}|p{7.93em}|l|l|l|p{4.855em}} \toprule
\textbf{事 件} & \textbf{地 點} & \multicolumn{1}{p{5.285em}|}{\textbf{馬太福音}} & \multicolumn{1}{p{5.57em}|}{\textbf{馬可福音}} & \multicolumn{1}{p{5.285em}|}{\textbf{路加福音}} & \textbf{約翰福音} \\  \midrule

\textcolor{red}{周日早}無花果樹被咒詛&從伯大尼出來& 21:18-19&11:12-14&&  \\\midrule

趕出做買賣/兌換銀錢/賣鴿子&耶路撒冷&21:12&  11:15-16 &19:45&  \\\midrule
不許人拿著器具從殿裡經過&耶路撒冷&&  &&  \\\midrule
责禱告的殿成為賊窩&耶路撒冷&  21:13  & 11:17&19:46&  \\\midrule
治好瞎子、瘸子&耶路撒冷&  21:14 &       &&  \\\midrule

祭司長文士惱怒小孩子赞美神&耶路撒冷&  21:15-16&&&  \\\midrule
应验從嬰孩吃奶口中完全讚美&耶路撒冷&21:15-16 &   && \\\midrule
祭司長文士等尊長要殺耶穌&耶路撒冷& & 11:18  &19:47-48&  \\\midrule
每天晚上耶穌出城去&伯大尼&21:17& 11:19&&  \\\midrule
 \end{tabular}%
\end{table}


\begin{itemize}
\item 早晨咒詛無花果樹
\begin{itemize}
馬太 21:18-19\\
早晨回城的時候，他餓了，
看見路旁有一棵無花果樹，就走到跟前，在樹上找不著什麼，不過有葉子，就對樹說：從今以後，你永不結果子。那無花果樹就立刻枯乾了。
\item 馬可 11:12-14\\
第二天，他們從伯大尼出來，耶穌餓了。
遠遠的看見一棵無花果樹，樹上有葉子，就往那裡去，或者在樹上可以找著什麼。到了樹下，竟找不著什麼，不過有葉子，因為不是收無花果的時候。
耶穌就對樹說：從今以後，永沒有人吃你的果子。他的門徒也聽見了。
\item 综合\\
第二天早晨,他們從伯大尼出來回城的時候，耶穌餓了。 遠遠地看見路旁有一棵無花果樹，就走到跟前，在樹上找不著什麼，不過有葉子，因為不是收無花果的時候。耶穌就對樹說：「從今以後，你永不結果子！永沒有人吃你的果子！」那無花果樹就立刻枯乾了, 他的門徒也聽見了。
\end{itemize}
\item 潔净聖殿
\begin{itemize}
\item 馬太 21:12-16\\
耶穌進了神的殿，趕出殿裡一切做買賣的人，推倒兌換銀錢之人的桌子，和賣鴿子之人的凳子；對他們說：經上記著說：我的殿必稱為禱告的殿，你們倒使他成為賊窩了。
在殿裡有瞎子、瘸子到耶穌跟前，他就治好了他們。祭司長和文士看見耶穌所行的奇事，又見小孩子在殿裡喊著說：和散那歸於大衛的子孫！就甚惱怒，
對他說：這些人所說的，你聽見了嗎？耶穌說：是的。經上說你從嬰孩和吃奶的口中完全了讚美的話，你們沒有念過嗎？
\item 馬可 11:15-17\\
他們來到耶路撒冷。耶穌進入聖殿，趕出殿裡做買賣的人，推倒兌換銀錢之人的桌子，和賣鴿子之人的凳子；
也不許人拿著器具從殿裡經過；
便教訓他們說：經上不是記著說：我的殿必稱為萬國禱告的殿嗎？你們倒使他成為賊窩了。
\item 路加 19:45-47a\\
耶稣进了殿，赶出里头作买卖的人， 对他们说，经上说，我的殿，必作祷告的殿。你们倒使它成为贼窝了。 耶稣天天在殿里教训人。 
%\item 综合\\
%耶穌來到耶路撒冷進入聖殿，趕出殿裡一切做買賣的人，推倒兌換銀錢之人的桌子和賣鴿子之人的凳子， 也不許人拿著器具從殿裡經過。 便教訓他們，說：「經上不是記著說『我的殿必稱為萬國禱告的殿』嗎？你們倒使它成為賊窩了！」在殿裡有瞎子、瘸子到耶穌跟前，他就治好了他們。
\end{itemize}
\item 祭司長等想滅耶穌\\
\begin{itemize}
\item 可11:18 \\祭司長和文士聽見這話，就想法子要除滅耶穌，卻又怕他，因為眾人都稀奇他的教訓。
\item 路19:47-48 \\耶穌天天在殿裡教訓人。祭司長和文士與百姓的尊長都想要殺他，但尋不出法子來，因為百姓都側耳聽他。
%\item 综合\\祭司長和文士聽見這話，就想法子要除滅耶穌，卻又怕他尋不出法子來，因為眾人都稀奇他的教訓側耳聽他。祭司長和文士看見耶穌所行的奇事，又見小孩子在殿裡喊著說「和散那歸於大衛的子孫」，就甚惱怒， 對他說：「這些人所說的，你聽見了嗎？」耶穌說：「是的。經上說『你從嬰孩和吃奶的口中完全了讚美』的話，你們沒有念過嗎？」
\end{itemize}

\item 夜宿伯大尼\\
\begin{itemize}
\item 馬太 21:17 - 於是離開他們，出城到伯大尼去，在那裡住宿。
\item 馬可 11:19 - 每天晚上，耶穌出城去。
\item 路加 21:37 - 耶穌每日在殿裡教訓人，每夜出城在一座山，名叫橄欖山住宿。
\end{itemize}
\end{itemize}





\subsubsection{耶稣的公开讲论（周一）}

\begin{table}[H]\centering
\begin{tabular}{p{14em}|p{7.93em}|l|l|l|p{4.855em}} \toprule
\textbf{事 件} & \textbf{地 點} & \multicolumn{1}{p{5.285em}|}{\textbf{馬太福音}} & \multicolumn{1}{p{5.57em}|}{\textbf{馬可福音}} & \multicolumn{1}{p{5.285em}|}{\textbf{路加福音}} & \textbf{約翰福音} \\  \midrule
\textcolor{red}{周一}早晨彼得见無花果樹枯乾&从伯大尼出来&21:20& 11:20-21&&  \\\midrule
对門徒论信心能移山&從伯大尼出來& 21:21-22&11:22-24&&  \\\midrule
禱告時饒恕人也被天父饒恕&從伯大尼出來& &11:25-26&&  \\\midrule
\end{tabular}%
\end{table}

\begin{itemize}
\item 馬太 21:20-22\\
門徒看見了，便希奇說：無花果樹怎麼立刻枯乾了呢？
耶穌回答說：我實在告訴你們，你們若有信心，不疑惑，不但能行無花果樹上所行的事，就是對這座山說：你挪開此地，投在海裡！也必成就。
你們禱告，無論求什麼，只要信，就必得著。
\item 馬可 11:20-26\\
早晨，他們從那裡經過，看見無花果樹連根都枯乾了。
彼得想起耶穌的話來，就對他說：拉比，請看！你所咒詛的無花果樹，已經枯乾了。
耶穌回答說：你們當信服神。
我實在告訴你們，無論何人對這座山說：你挪開此地，投在海裡！他若心裡不疑惑，只信他所說的必成，就必給他成了。所以我告訴你們，凡你們禱告祈求的，無論是什麼，只要信是得著的，就必得著。
你們站著禱告的時候，若想起有人得罪你們，就當饒恕他，好叫你們在天上的父也饒恕你們的過犯。
你們若不饒恕人，你們在天上的父也不饒恕你們的過犯。（有古卷無此節）
\item 综合\\
早晨，他們回城的時候從那裡經過，看見無花果樹連根都枯乾了。 彼得想起耶穌的話來，就對他說：「拉比，請看，你所咒詛的無花果樹已經枯乾了！」耶穌回答說：「你們當信服神。我實在告訴你們：你們若有信心，不疑惑，不但能行無花果樹上所行的事，無論何人就是對這座山說：『你挪開此地，投在海裡！』，他若心裡不疑惑，只信他所說的必成，就必給他成了。所以我告訴你們，凡你們禱告祈求的，無論是什麼，只要信是得著的，就必得著。 你們站著禱告的時候，若想起有人得罪你們，就當饒恕他，好叫你們在天上的父也饒恕你們的過犯。 你們若不饒恕人，你們在天上的父也不饒恕你們的過犯。」
\end{itemize}


\begin{table}[H]\centering
\begin{tabular}{p{14em}|p{7.93em}|l|l|l|p{4.855em}} \toprule
\textbf{事 件} & \textbf{地 點} & \multicolumn{1}{p{5.285em}|}{\textbf{馬太福音}} & \multicolumn{1}{p{5.57em}|}{\textbf{馬可福音}} & \multicolumn{1}{p{5.285em}|}{\textbf{路加福音}} & \textbf{約翰福音} \\  \midrule
祭司長文士並長老辯權柄&殿 & 21:23-27 &11:27-33&20:1-8&  \\\midrule
兩個兒子葡萄園做工比喻&殿 & 21:28-32 &&&  \\\midrule
兇惡園戶租葡萄園的比喻&殿 &21:33-44 & 12:1-11 &20:9-18&  \\\midrule
祭司長法利賽人要捉拿耶稣&殿 &21:45-46&  12:12     &&  \\\midrule

王為兒子娶親的筵席&殿 &22:1-14&      &&  \\\midrule

文士祭司長法利賽人派人問納稅& 殿   &  22:15-22     &12:13-17  &20:19-26&  \\\midrule
撒都該人文士問一妻嫁七兄& 殿 & 22:23-33 & 12:18-27 &20:27-39&  \\\midrule

法利賽人律法師问最大誡命&殿 &22:34-40&  12:28-34  &&  \\\midrule
耶穌問基督是大衛子孫&   殿    &22:41-46  &  12:35-37  &20:40-44&  \\\midrule
耶穌戒門徒眾人防文士假冒伪善&殿&23:1-12   &&&  \\\midrule

法利賽人的七禍& 殿&23:13-36   &12:38-40&20:45-47&  \\\midrule

嘆耶路撒冷常殺先知成為荒場& 殿&23:37-39 & &       &  \\\midrule

讚窮寡婦的奉獻&耶穌對銀庫坐&       &12:41-44&21:1-4&  \\\midrule

%耶穌向百姓末了的公开講論&&       &&&  \\\midrule
麥子落在地裡死了結出許多子粒&&       &&&12:23-24\\\midrule
惜命失命,恨命保守生命到永生&&       &&&12:25 \\\midrule
我在哪裡服侍我的人也要在哪裡&&       &&&12:26  \\\midrule
父子对话:榮耀父名還要再榮耀&&       &&&12:27-28  \\\midrule
眾人聽見天上聲音的反应&&       &&&12:29\\\midrule
世界受審判世界的王要被趕出去&&       &&&12:30-31\\\midrule
從地上被舉起吸引萬人來歸我&&       &&&12:32-33\\\midrule
對眾人說趁著有光成為光明之子&&       &&&12:34-36\\\midrule
眾人不信&&       &&&12:34-50\\\midrule
\end{tabular}%
\end{table}



\begin{itemize}
\item 質問耶穌權柄 -太 21:23-27/可11:27-33/路 20:1-8
\begin{itemize}
\item{馬太 21:23-27}
耶穌進了殿，正教訓人的時候，祭司長和民間的長老來問他說：你仗著什麼權柄做這些事？給你這權柄的是誰呢？耶穌回答說：我也要問你們一句話，你們若告訴我，我就告訴你們我仗著什麼權柄做這些事。約翰的洗禮是從哪裡來的？是從天上來的？是從人間來的呢？他們彼此商議說：我們若說從天上來，他必對我們說：這樣，你們為什麼不信他呢？若說從人間來，我們又怕百姓，因為他們都以約翰為先知。於是回答耶穌說：我們不知道。耶穌說：我也不告訴你們我仗著什麼權柄做這些事。
\item{馬可 11:27-33 }
他們又來到耶路撒冷。耶穌在殿裡行走的時候，祭司長和文士並長老進前來，
問他說：你仗著什麼權柄做這些事？給你這權柄的是誰呢？耶穌對他們說：我要問你們一句話，你們回答我，我就告訴你們我仗著什麼權柄做這些事。約翰的洗禮是從天上來的？是從人間來的呢？你們可以回答我。
他們彼此商議說：我們若說從天上來，他必說：這樣，你們為什麼不信他呢？
若說從人間來，卻又怕百姓，因為眾人真以約翰為先知。
於是回答耶穌說：我們不知道。耶穌說：我也不告訴你們我仗著什麼權柄做這些事。
\item{路加 20:1-8}
有一天，耶穌在殿裡教訓百姓，講福音的時候，祭司長和文士並長老上前來，
問他說：你告訴我們，你仗著什麼權柄做這些事？給你這權柄的是誰呢？
耶穌回答說：我也要問你們一句話，你們且告訴我。
約翰的洗禮是從天上來的？是從人間來的呢？
他們彼此商議說：我們若說從天上來，他必說：你們為什麼不信他呢？
若說從人間來，百姓都要用石頭打死我們，因為他們信約翰是先知。
於是回答說：不知道是從哪裡來的。
耶穌說：我也不告訴你們，我仗著什麼權柄做這些事。
\item{综合}
耶穌和門徒又來到耶路撒冷進了殿，正教訓人講福音的時候(在殿裡行走的時候)，祭司長和文士並長老進前來， 問他說：「你仗著什麼權柄做這些事？給你這權柄的是誰呢？」 耶穌對他們說：「我要問你們一句話，你們回答我，我就告訴你們我仗著什麼權柄做這些事。 約翰的洗禮是從天上來的，是從人間來的呢？你們可以回答我。」 他們彼此商議說：「我們若說從天上來，他必說：『這樣，你們為什麼不信他呢？』 若說從人間來，卻又怕百姓，百姓都要用石頭打死我們，因為眾人真以約翰為先知。」於是回答耶穌說：「我們不知道是從哪裡來的。」耶穌說：「我也不告訴你們我仗著什麼權柄做這些事。」
\end{itemize}
\item{兩個兒子去葡萄園做工的比喻 - 馬太 21:28-32}\\
%\paragraph{馬太 21:28-32}
又說：一個人有兩個兒子。他來對大兒子說：我兒，你今天到葡萄園裡去做工。
他回答說：我不去，以後自己懊悔，就去了。
又來對小兒子也是這樣說。他回答說：父啊，我去，他卻不去。
你們想，這兩個兒子是哪一個遵行父命呢？他們說：大兒子。耶穌說：我實在告訴你們，稅吏和娼妓倒比你們先進神的國。
因為約翰遵著義路到你們這裡來，你們卻不信他；稅吏和娼妓倒信他。你們看見了，後來還是不懊悔去信他。
\item{凶惡葡萄園戶-太21:33-44/可12:1-11/路20:9-18}
\begin{itemize}
\item{馬太 21:33-44}
你們再聽一個比喻：有個家主栽了一個葡萄園，周圍圈上籬笆，裡面挖了一個壓酒池，蓋了一座樓，租給園戶，就往外國去了。
收果子的時候近了，就打發僕人到園戶那裡去收果子。
園戶拿住僕人，打了一個，殺了一個，用石頭打死一個。
主人又打發別的僕人去，比先前更多；園戶還是照樣待他們。
後來打發他的兒子到他們那裡去，意思說：他們必尊敬我的兒子。
不料，園戶看見他兒子，就彼此說：這是承受產業的。來吧，我們殺他，佔他的產業！
他們就拿住他，推出葡萄園外，殺了。
園主來的時候要怎樣處治這些園戶呢？
他們說：要下毒手除滅那些惡人，將葡萄園另租給那按著時候交果子的園戶。
耶穌說：經上寫著：匠人所棄的石頭已作了房角的頭塊石頭。這是主所做的，在我們眼中看為希奇。這經你們沒有念過嗎？
所以我告訴你們，神的國必從你們奪去，賜給那能結果子的百姓。
誰掉在這石頭上，必要跌碎；這石頭掉在誰的身上，就要把誰砸得稀爛。
\item{馬可 12:1-11}
耶穌就用比喻對他們說：有人栽了一個葡萄園，周圍圈上籬笆，挖了一個壓酒池，蓋了一座樓，租給園戶，就往外國去了。
到了時候，打發一個僕人到園戶那裡，要從園戶收葡萄園的果子。
園戶拿住他，打了他，叫他空手回去。
再打發一個僕人到他們那裡。他們打傷他的頭，並且凌辱他。
又打發一個僕人去，他們就殺了他。後又打發好些僕人去，有被他們打的，有被他們殺的。
園主還有一位是他的愛子，末後又打發他去，意思說：他們必尊敬我的兒子。
不料，那些園戶彼此說：這是承受產業的。來吧，我們殺他，產業就歸我們了！
於是拿住他，殺了他，把他丟在園外。
這樣，葡萄園的主人要怎麼辦呢？他要來除滅那些園戶，將葡萄園轉給別人。
經上寫著說：匠人所棄的石頭，已作了房角的頭塊石頭。
這是主所做的，在我們眼中看為希奇。這經你們沒有念過嗎？
\item{路加 20:9-18}
耶穌就設比喻對百姓說：有人栽了一個葡萄園，租給園戶，就往外國去住了許久。
到了時候，打發一個僕人到園戶那裡去，叫他們把園中當納的果子交給他；園戶竟打了他，叫他空手回去。
又打發一個僕人去，他們也打了他，並且凌辱他，叫他空手回去。
又打發第三個僕人去，他們也打傷了他，把他推出去了。
園主說：我怎麼辦呢？我要打發我的愛子去，或者他們尊敬他。
不料，園戶看見他，就彼此商量說：這是承受產業的，我們殺他吧，使產業歸於我們！
於是把他推出葡萄園外，殺了。這樣，葡萄園的主人要怎樣處治他們呢？
他要來除滅這些園戶，將葡萄園轉給別人。聽見的人說：這是萬不可的！
耶穌看著他們說：經上記著：匠人所棄的石頭已作了房角的頭塊石頭。這是什麼意思呢？
凡掉在那石頭上的，必要跌碎；那石頭掉在誰的身上，就要把誰砸得稀爛。
文士和祭司長看出這比喻是指著他們說的，當時就想要下手拿他，只是懼怕百姓。
\item{综合}
耶穌就用比喻對他們說：「有人栽了一個葡萄園，周圍圈上籬笆，挖了一個壓酒池，蓋了一座樓，租給園戶，就往外國去住了許久。 到了時候，打發一個僕人到園戶那裡，要從園戶收葡萄園的果子, 叫他們把園中當納的果子交給他。 園戶拿住他，打了他，叫他空手回去。  再打發一個僕人到他們那裡，他們打傷他的頭，並且凌辱他，叫他空手回去。 又打發三個個僕人去，他們就打傷了他，把他推出去殺了他。後又打發好些僕人去，比先前更多，園戶還是照樣待他們,有被他們打的，有被他們殺的。園主說：『我怎麼辦呢？我要打發我的愛子去，或者他們尊敬他。』還有一位是他的愛子，末後又打發他去，意思說：『他們必尊敬我的兒子。』 不料那些園戶彼此說：『這是承受產業的。來吧，我們殺他，占他的產業！產業就歸我們了！』 於是拿住他，殺了他，把他丟在園外。 這樣，葡萄園的主人要怎麼辦呢？他們說：「要下毒手除滅那些惡人，將葡萄園另租給那按著時候交果子的園戶。」耶穌說：他要來除滅那些園戶，將葡萄園轉給別人。  聽見的人說：「這是萬不可的！」耶穌看著他們說：經上寫著說：『匠人所棄的石頭，已做了房角的頭塊石頭。這是主所做的，在我們眼中看為稀奇。』這是什麼意思呢？凡掉在那石頭上的，必要跌碎；那石頭掉在誰的身上，就要把誰砸得稀爛。」這經你們沒有念過嗎？所以我告訴你們：神的國必從你們奪去，賜給那能結果子的百姓。」他們看出這比喻是指著他們說的，就想要捉拿他，只是懼怕百姓，因為眾人以他為先知,於是離開他走了。
\end{itemize}
\item{王的兒子娶親筵席的比喻 - 馬太22:1-14}\\
耶穌又用比喻對他們說：
天國好比一個王為他兒子擺設娶親的筵席，
就打發僕人去，請那些被召的人來赴席，他們卻不肯來。
王又打發別的僕人，說：你們告訴那被召的人，我的筵席已經預備好了，牛和肥畜已經宰了，各樣都齊備，請你們來赴席。
那些人不理就走了；一個到自己田裡去；一個做買賣去；
其餘的拿住僕人，凌辱他們，把他們殺了。
王就大怒，發兵除滅那些兇手，燒毀他們的城。
於是對僕人說：喜筵已經齊備，只是所召的人不配。
所以你們要往岔路口上去，凡遇見的，都召來赴席。
那些僕人就出去，到大路上，凡遇見的，不論善惡都召聚了來，筵席上就坐滿了客。
王進來觀看賓客，見那裡有一個沒有穿禮服的，
就對他說：朋友，你到這裡來怎麼不穿禮服呢？那人無言可答。
於是王對使喚的人說：捆起他的手腳來，把他丟在外邊的黑暗裡；在那裡必要哀哭切齒了。
因為被召的人多，選上的人少。

\item{納稅給凱撒-太22:15-22/可12:13-17/路20:19-26}
\begin{itemize}
\item{馬太 22:15-22}
當時，法利賽人出去商議，怎樣就著耶穌的話陷害他，
就打發他們的門徒同希律黨的人去見耶穌，說：夫子，我們知道你是誠實人，並且誠誠實實傳神的道，什麼人你都不徇情面，因為你不看人的外貌。
請告訴我們，你的意見如何？納稅給該撒可以不可以？
耶穌看出他們的惡意，就說：假冒為善的人哪，為什麼試探我？
拿一個上稅的錢給我看！他們就拿一個銀錢來給他。
耶穌說：這像和這號是誰的？
他們說：是該撒的。耶穌說：這樣，該撒的物當歸給該撒；神的物當歸給神。
他們聽見就希奇，離開他走了。
\item{馬可 12:13-17}
後來，他們打發幾個法利賽人和幾個希律黨的人到耶穌那裡，要就著他的話陷害他。
他們來了，就對他說：夫子，我們知道你是誠實的，什麼人你都不徇情面；因為你不看人的外貌，乃是誠誠實實傳神的道。納稅給該撒可以不可以？
我們該納不該納？耶穌知道他們的假意，就對他們說：你們為什麼試探我？拿一個銀錢來給我看！
他們就拿了來。耶穌說：這像和這號是誰的？他們說：是該撒的。
耶穌說：該撒的物當歸給該撒，神的物當歸給神。他們就很希奇他。
\item{路加20:19-26}
於是窺探耶穌，打發奸細裝作好人，要在他的話上得把柄，好將他交在巡撫的政權之下。
奸細就問耶穌說：夫子，我們曉得你所講所傳都是正道，也不取人的外貌，乃是誠誠實實傳神的道。
我們納稅給該撒，可以不可以？
耶穌看出他們的詭詐，就對他們說：
拿一個銀錢來給我看。這像和這號是誰的？他們說：是該撒的。
耶穌說：這樣，該撒的物當歸給該撒，神的物當歸給神。
他們當著百姓，在這話上得不著把柄，又希奇他的應對，就閉口無言了。
\item{综合}
法利賽人出去商議，怎樣就著耶穌的話陷害他，就打發他們的門徒同希律黨的人去窺探耶穌，打發奸細裝作好人，要在他的話上得把柄，好將他交在巡撫的政權之下。奸細就問耶穌說：「夫子，我們知道你是誠實人，並且誠誠實實傳神的道，什麼人你都不徇情面，因為你不看人的外貌。請告訴我們你的意見如何：納稅給愷撒可以不可以? 我們該納不該納？」耶穌知道他們的假意，就對他們說：「假冒為善的人哪，你們為什麼試探我？拿一個上稅的銀錢來給我看。」他們就拿了來。耶穌說：「這像和這號是誰的？」他們說：「是愷撒的。」耶穌說：「愷撒的物當歸給愷撒，神的物當歸給神。」他們當著百姓，在這話上得不著把柄，又稀奇他的應對，就閉口無言了, 離開他走了。
\end{itemize}

\item{撒都該人問復活-太22:23-33/可12:18-27/路20:27-39}
\begin{itemize}
\item{馬太 22:23-33}
撒都該人常說沒有復活的事。那天，他們來問耶穌說：
夫子，摩西說：人若死了，沒有孩子，他兄弟當娶他的妻，為哥哥生子立後。
從前，在我們這裡有弟兄七人，第一個娶了妻，死了，沒有孩子，撇下妻子給兄弟。
第二、第三、直到第七個，都是如此。
末後，婦人也死了。
這樣，當復活的時候，他是七個人中哪一個的妻子呢？因為他們都娶過他。
耶穌回答說：你們錯了；因為不明白聖經，也不曉得神的大能。
當復活的時候，人也不娶也不嫁，乃像天上的使者一樣。
論到死人復活，神在經上向你們所說的，你們沒有念過嗎？
他說：我是亞伯拉罕的神，以撒的神，雅各的神。神不是死人的神，乃是活人的神。
眾人聽見這話，就希奇他的教訓。
\item{馬可 12:18-27}
撒都該人常說沒有復活的事。他們來問耶穌說：
夫子，摩西為我們寫著說：人若死了，撇下妻子，沒有孩子，他兄弟當娶他的妻，為哥哥生子立後。
有弟兄七人，第一個娶了妻，死了，沒有留下孩子。
第二個娶了他，也死了，沒有留下孩子。第三個也是這樣。
那七個人都沒有留下孩子；末了，那婦人也死了。
當復活的時候，他是哪一個的妻子呢？因為他們七個人都娶過他。
耶穌說：你們所以錯了，豈不是因為不明白聖經，不曉得神的大能嗎？
人從死裡復活，也不娶也不嫁，乃像天上的使者一樣。
論到死人復活，你們沒有念過摩西的書荊棘篇上所載的嗎？神對摩西說：我是亞伯拉罕的神，以撒的神，雅各的神。
神不是死人的神，乃是活人的神。你們是大錯了。
\item{路加 20:27-39}
撒都該人常說沒有復活的事。有幾個來問耶穌說：
夫子！摩西為我們寫著說：人若有妻無子就死了，他兄弟當娶他的妻，為哥哥生子立後。
有弟兄七人，第一個娶了妻，沒有孩子死了；
第二個、第三個也娶過他；
那七個人都娶過他，沒有留下孩子就死了。
後來婦人也死了。
這樣，當復活的時候，他是哪一個的妻子呢？因為他們七個人都娶過他。
耶穌說：這世界的人有娶有嫁；
惟有算為配得那世界，與從死裡復活的人也不娶也不嫁；
因為他們不能再死，和天使一樣；既是復活的人，就為神的兒子。
至於死人復活，摩西在荊棘篇上，稱主是亞伯拉罕的神，以撒的神，雅各的神，就指示明白了。
神原不是死人的神，乃是活人的神；因為在他那裡（那裡：或作看來），人都是活的。有幾個文士說：夫子！你說得好。

\item{综合}
撒都該人常說沒有復活的事。有幾個來問耶穌，說：「夫子，摩西為我們寫著說：『人若死了，撇下妻子，沒有孩子，他兄弟當娶他的妻，為哥哥生子立後。』有弟兄七人，第一個娶了妻，死了，沒有留下孩子。第二個娶了她，也死了，沒有留下孩子。第三個也是這樣。那七個人都沒有留下孩子。末了，那婦人也死了。當復活的時候，她是哪一個的妻子呢？因為他們七個人都娶過她。」耶穌說：「你們所以錯了，豈不是因為不明白聖經，不曉得神的大能嗎？這世界的人有娶有嫁，唯有算為配得那世界與從死裡復活的人，也不娶也不嫁，乃像天上的使者一樣。因為他們不能再死，和天使一樣；既是復活的人，就為神的兒子。論到死人復活，你們沒有念過摩西的書『荊棘篇』上所載的嗎？神對摩西說：『我是亞伯拉罕的神、以撒的神、雅各的神』就指示明白了。神不是死人的神，乃是活人的神。因為在他那裡人都是活的。你們是大錯了！」眾人聽見這話，就稀奇他的教訓。
\end{itemize}
\item{論最大的誡命-太22:34-40/可12:28-34}
\begin{itemize}
\item{馬太 22:34-40}
法利賽人聽見耶穌堵住了撒都該人的口，他們就聚集。
內中有一個人是律法師，要試探耶穌，就問他說：
夫子，律法上的誡命，哪一條是最大的呢？
耶穌對他說：你要盡心、盡性、盡意愛主─你的神。
這是誡命中的第一，且是最大的。
其次也相倣，就是要愛人如己。
這兩條誡命是律法和先知一切道理的總綱。
\item{馬可 12:28-34}
有一個文士來，聽見他們辯論，曉得耶穌回答的好，就問他說：誡命中哪是第一要緊的呢？
耶穌回答說：第一要緊的就是說：以色列啊，你要聽，主─我們神是獨一的主。
你要盡心、盡性、盡意、盡力愛主─你的神。
其次就是說：要愛人如己。再沒有比這兩條誡命更大的了。
那文士對耶穌說：夫子說，神是一位，實在不錯；除了他以外，再沒有別的神；
並且盡心、盡智、盡力愛他，又愛人如己，就比一切燔祭和各樣祭祀好的多。
耶穌見他回答的有智慧，就對他說：你離神的國不遠了。從此以後，沒有人敢再問他什麼。
\item{综合}
法利賽人聽見耶穌堵住了撒都該人的口，他們就聚集。內中有一個人是律法師(文士)來，聽見他們辯論，曉得耶穌回答得好，要試探耶穌, 就問他說：「誡命中哪是第一要緊(最大)的呢？」耶穌回答說：「第一要緊的且是最大的就是說：『以色列啊，你要聽，主我們神是獨一的主。你要盡心、盡性、盡意、盡力愛主你的神。』其次也相仿就是說：『要愛人如己。』再沒有比這兩條誡命更大的了,這兩條誡命是律法和先知一切道理的總綱。」那文士對耶穌說：「夫子說神是一位，實在不錯！除了他以外，再沒有別的神。並且盡心、盡智、盡力愛他，又愛人如己，就比一切燔祭和各樣祭祀好得多。」耶穌見他回答得有智慧，就對他說：「你離神的國不遠了。」從此以後，沒有人敢再問他什麼。有幾個文士說：「夫子，你說得好！」以後他們不敢再問他什麼。
\end{itemize}

\item{論基督與大衛關係-太22:41-46/可12:35-37/路20:40-44}
\begin{itemize}
\item{馬太 22:41-46}
法利賽人聚集的時候，耶穌問他們說：
論到基督，你們的意見如何？他是誰的子孫呢？他們回答說：是大衛的子孫。
耶穌說：這樣，大衛被聖靈感動，怎麼還稱他為主，說：
主對我主說：你坐在我的右邊，等我把你仇敵放在你的腳下。
大衛既稱他為主，他怎麼又是大衛的子孫呢？
他們沒有一個人能回答一言。從那日以後，也沒有人敢再問他什麼。
\item{馬可 12:35-37}
耶穌在殿裡教訓人，就問他們說：文士怎麼說基督是大衛的子孫呢？
大衛被聖靈感動，說：主對我主說，你坐在我的右邊，等我使你仇敵作你的腳凳。
大衛既自己稱他為主，他怎麼又是大衛的子孫呢？眾人都喜歡聽他。
\item{路加 20:40-44}
以後他們不敢再問他什麼。耶穌對他們說：人怎麼說基督是大衛的子孫呢？詩篇上大衛自己說：主對我主說：你坐在我的右邊，等我使你仇敵作你的腳凳。大衛既稱他為主，他怎麼又是大衛的子孫呢？
\item{综合}
耶穌在殿裡教訓人，法利賽人聚集的時候，就問他們說：「論到基督，你們的意見如何？他是誰的子孫呢？」他們回答說：「是大衛的子孫。」耶穌說：「文士怎麼說基督是大衛的子孫呢？這樣，大衛被聖靈感動，還稱他為主說：『主對我主說：「你坐在我的右邊，等我使你仇敵做你的腳凳。」』大衛既自己稱他為主，他怎麼又是大衛的子孫呢？」他們沒有一個人能回答一言。從那日以後，也沒有人敢再問他什麼。眾人都喜歡聽他。
\end{itemize}
\item{責備法利賽人文士-太23:1-12}
那時，耶穌對眾人和門徒講論，
說：文士和法利賽人坐在摩西的位上，
凡他們所吩咐你們的，你們都要謹守遵行；但不要效法他們的行為；因為他們能說，不能行。
他們把難擔的重擔捆起來，擱在人的肩上，但自己一個指頭也不肯動。
他們一切所做的事都是要叫人看見，所以將佩戴的經文做寬了，衣裳的繸子做長了，
喜愛筵席上的首座，會堂裡的高位，
又喜愛人在街市上問他安，稱呼他拉比（拉比就是夫子）。
但你們不要受拉比的稱呼，因為只有一位是你們的夫子；你們都是弟兄。
也不要稱呼地上的人為父，因為只有一位是你們的父，就是在天上的父。
也不要受師尊的稱呼，因為只有一位是你們的師尊，就是基督。
你們中間誰為大，誰就要作你們的用人。
凡自高的，必降為卑；自卑的，必升為高。

\item{責備法利賽人文士-太23:13-36/可12:38-40/路20:45-47}
\begin{itemize}
\item{馬太 23:1-36}
你們這假冒為善的文士和法利賽人有禍了！因為你們正當人前，把天國的門關了，自己不進去，正要進去的人，你們也不容他們進去。（有古卷在此有
你們這假冒為善的文士和法利賽人有禍了！因為你們侵吞寡婦的家產，假意做很長的禱告，所以要受更重的刑罰。）
你們這假冒為善的文士和法利賽人有禍了！因為你們走遍洋海陸地，勾引一個人入教，既入了教，卻使他作地獄之子，比你們還加倍。
你們這瞎眼領路的有禍了！你們說：凡指著殿起誓的，這算不得什麼；只是凡指著殿中金子起誓的，他就該謹守。
你們這無知瞎眼的人哪，什麼是大的？是金子呢？還是叫金子成聖的殿呢？
你們又說：凡指著壇起誓的，這算不得什麼；只是凡指著壇上禮物起誓的，他就該謹守。
你們這瞎眼的人哪，什麼是大的？是禮物呢？還是叫禮物成聖的壇呢？
所以，人指著壇起誓，就是指著壇和壇上一切所有的起誓；
人指著殿起誓，就是指著殿和那住在殿裡的起誓；
人指著天起誓，就是指著神的寶座和那坐在上面的起誓。
你們這假冒為善的文士和法利賽人有禍了！因為你們將薄荷、茴香、芹菜，獻上十分之一，那律法上更重的事，就是公義、憐憫、信實，反倒不行了。這更重的是你們當行的；那也是不可不行的。
你們這瞎眼領路的，蠓蟲你們就濾出來，駱駝你們倒吞下去。
你們這假冒為善的文士和法利賽人有禍了！因為你們洗淨杯盤的外面，裡面卻盛滿了勒索和放蕩。
你這瞎眼的法利賽人，先洗淨杯盤的裡面，好叫外面也乾淨了。
你們這假冒為善的文士和法利賽人有禍了！因為你們好像粉飾的墳墓，外面好看，裡面卻裝滿了死人的骨頭和一切的污穢。
你們也是如此，在人前，外面顯出公義來，裡面卻裝滿了假善和不法的事。
你們這假冒為善的文士和法利賽人有禍了！因為你們建造先知的墳，修飾義人的墓，說：
若是我們在我們祖宗的時候，必不和他們同流先知的血。
這就是你們自己證明是殺害先知者的子孫了。
你們去充滿你們祖宗的惡貫吧！
你們這些蛇類、毒蛇之種啊，怎能逃脫地獄的刑罰呢？
所以我差遣先知和智慧人並文士到你們這裡來，有的你們要殺害，要釘十字架；有的你們要在會堂裡鞭打，從這城追逼到那城，
叫世上所流義人的血都歸到你們身上，從義人亞伯的血起，直到你們在殿和壇中間所殺的巴拉加的兒子撒迦利亞的血為止。
我實在告訴你們，這一切的罪都要歸到這世代了。

\item{馬可 12:38-40}
耶穌在教訓之間，說：你們要防備文士；他們好穿長衣遊行，喜愛人在街市上問他們的安，
又喜愛會堂裡的高位，筵席上的首座。
他們侵吞寡婦的家產，假意做很長的禱告。這些人要受更重的刑罰！
\item{路加 20:45-47}
眾百姓聽的時候，耶穌對門徒說：
你們要防備文士。他們好穿長衣遊行，喜愛人在街市上問他們安，又喜愛會堂裡的高位，筵席上的首座；
他們侵吞寡婦的家產，假意作很長的禱告。這些人要受更重的刑罰！
\item{综合}
那時，耶穌對眾人和門徒講論，說：「文士和法利賽人坐在摩西的位上，凡他們所吩咐你們的，你們都要謹守遵行，但不要效法他們的行為，因為他們能說不能行。他們把難擔的重擔捆起來，擱在人的肩上，但自己一個指頭也不肯動。他們一切所做的事都是要叫人看見，所以將佩戴的經文做寬了，衣裳的穗子做長了；喜愛筵席上的首座、會堂裡的高位，又喜愛人在街市上問他安，稱呼他『拉比』 。但你們不要受拉比的稱呼，因為只有一位是你們的夫子，你們都是弟兄。也不要稱呼地上的人為父，因為只有一位是你們的父，就是在天上的父。也不要受師尊的稱呼，因為只有一位是你們的師尊，就是基督。你們中間誰為大，誰就要做你們的用人。凡自高的，必降為卑；自卑的，必升為高。
假冒為善的文士和法利賽人	因為你們正當人前把天國的門關了，自己不進去，正要進去的人你們也不容他們進去。（太23:13）
假冒為善的文士和法利賽人	侵吞寡婦的家產，假意做很長的禱告，所以要受更重的刑罰。  （太23：14）
假冒為善的文士和法利賽人	因為你們走遍洋海陸地，勾引一個人入教，既入了教，卻使他做地獄之子，比你們還加倍。（太23:15）
瞎眼領路的	凡指著殿起誓的，這算不得什麼；凡指著壇起誓的，這算不得什麼。（太23：17-22）
假冒為善的文士和法利賽人	因為你們將薄荷、茴香、芹菜獻上十分之一，那律法上更重的事，就是公義、憐憫、信實反倒不行了。（太23:23,24）
假冒為善的文士和法利賽人	因為你們洗淨杯盤的外面，裡面卻盛滿了勒索和放蕩。（太23：25,26）
假冒為善的文士和法利賽人	在人前，外面顯出公義來，裡面卻裝滿了假善和不法的事。（太23:27,28）
假冒為善的文士和法利賽人	所以我差遣先知和智慧人並文士到你們這裡來，有的你們要殺害，要釘十字架，有的你們要在會堂裡鞭打，從這城追逼到那城， 35 叫世上所流義人的血都歸到你們身上。（太23:29-33）
文士	他們好穿長衣遊行，喜愛人在街市上問他們的安，又喜愛會堂裡的高位、筵席上的首座。他們侵吞寡婦的家產，假意作很長的禱告。(路20：46，47)
\end{itemize}
\item{主稱讚寡婦的捐資-馬可 12:41-44/路加 21:1-4}
\begin{itemize}
\item{馬可 12:41-44}
耶穌對銀庫坐著，看眾人怎樣投錢入庫。有好些財主往裡投了若干的錢。
有一個窮寡婦來，往裡投了兩個小錢，就是一個大錢。
耶穌叫門徒來，說：我實在告訴你們，這窮寡婦投入庫裡的，比眾人所投的更多。
因為，他們都是自己有餘，拿出來投在裡頭；但這寡婦是自己不足，把他一切養生的都投上了。
\item{路加 21:1-4}
耶穌擡頭觀看，見財主把捐項投在庫裡，又見一個窮寡婦投了兩個小錢，就說：我實在告訴你們，這窮寡婦所投的比眾人還多；因為眾人都是自己有餘，拿出來投在捐項裡，但這寡婦是自己不足，把他一切養生的都投上了。
\item{综合}
耶穌對銀庫坐著抬頭觀看，看眾人怎樣投錢入庫。有好些財主往裡投了若干的錢。有一個窮寡婦來，往裡投了兩個小錢，就是一個大錢。耶穌叫門徒來，說：「我實在告訴你們：這窮寡婦投入庫裡的比眾人所投的更多，因為他們都是自己有餘，拿出來投在裡頭，但這寡婦是自己不足，把她一切養生的都投上了。」
\end{itemize}
%\subsection{希利尼人來見耶穌  (約12:20-22）}
%那時，上來過節禮拜的人中，有幾個希臘人。他們來見加利利伯賽大的腓力，求他說：「先生，我們願意見耶穌。」腓力去告訴安得烈，安得烈同腓力去告訴耶穌。
\item{最後的公開講論（約12:23-50）}
耶穌說：「人子得榮耀的時候到了。我實實在在地告訴你們：一粒麥子不落在地裡死了，仍舊是一粒；若是死了，就結出許多子粒來。愛惜自己生命的，就失喪生命；在這世上恨惡自己生命的，就要保守生命到永生。若有人服侍我，就當跟從我；我在哪裡，服侍我的人也要在哪裡。若有人服侍我，我父必尊重他。我現在心裡憂愁，我說什麼才好呢？父啊，救我脫離這時候！但我原是為這時候來的。父啊，願你榮耀你的名！」當時就有聲音從天上來說：「我已經榮耀了我的名，還要再榮耀。」站在旁邊的眾人聽見，就說：「打雷了。」還有人說：「有天使對他說話。」耶穌說：「這聲音不是為我，是為你們來的。現在這世界受審判，這世界的王要被趕出去。我若從地上被舉起來，就要吸引萬人來歸我。」耶穌這話原是指著自己將要怎樣死說的。眾人回答說：「我們聽見律法上有話說基督是永存的，你怎麼說人子必須被舉起來呢？這人子是誰呢？」耶穌對他們說：「光在你們中間還有不多的時候，應當趁著有光行走，免得黑暗臨到你們。那在黑暗裡行走的，不知道往何處去。你們應當趁著有光，信從這光，使你們成為光明之子。」
耶穌說了這話，就離開他們隱藏了。他雖然在他們面前行了許多神蹟，他們還是不信他。這是要應驗先知以賽亞的話說：「主啊，我們所傳的有誰信呢？主的膀臂向誰顯露呢？」他們所以不能信，因為以賽亞又說：「主叫他們瞎了眼，硬了心，免得他們眼睛看見，心裡明白，回轉過來，我就醫治他們。」以賽亞因為看見他的榮耀，就指著他說這話。雖然如此，官長中卻有好些信他的，只因法利賽人的緣故，就不承認，恐怕被趕出會堂。這是因他們愛人的榮耀過於愛神的榮耀。
耶穌大聲說：「信我的，不是信我，乃是信那差我來的。人看見我，就是看見那差我來的。我到世上來，乃是光，叫凡信我的不住在黑暗裡。若有人聽見我的話不遵守，我不審判他。我來本不是要審判世界，乃是要拯救世界。棄絕我、不領受我話的人，有審判他的，就是我所講的道在末日要審判他。因為我沒有憑著自己講，唯有差我來的父已經給我命令，叫我說什麼、講什麼。我也知道他的命令就是永生。故此，我所講的話正是照著父對我所說的。」
\end{itemize}


%\subsubsection{主降臨和世界末了預兆}
\begin{table}[H]\centering
\begin{tabular}{p{14em}|p{7.93em}|l|l|l|p{4.855em}}\toprule
 \textbf{事 件} & \textbf{地 點} & \multicolumn{1}{p{5.285em}|}{\textbf{馬太福音}} & \multicolumn{1}{p{5.57em}|}{\textbf{馬可福音}} & \multicolumn{1}{p{5.285em}|}{\textbf{路加福音}} & \textbf{約翰福音} \\  \midrule
預言聖殿被毀&走出聖殿時&24:1-2&13:1-2&21:5-6&  \\\midrule
彼得等暗问降臨和世界末了預兆&雅各約翰安得烈&24:3&13:3-4&21:7-8&  \\\midrule
冒主名來的“基督”迷惑多人&橄欖山上坐著&24:4-5&13:5-6&21:7-8&  \\\midrule
打仗和擾亂末期没立時到& &24:6&13:7 &21:9&  \\\midrule

民要攻打民，國要攻打國&&24:7a&13:8a&21:10&  \\\midrule
多處有饑荒地震災難起頭&&24:7b-8&13:8b&21:10&  \\\midrule

你們入監遭逼迫被殺被萬民恨惡&       &24:9&13:9 &21:12-13&  \\\midrule
多人跌倒彼此陷害、彼此恨惡&       &24:10&13:12-13a&&  \\\midrule
好些假先知起來迷惑多人&       &24:11&       &&  \\\midrule
因不法事增多，多人愛心漸冷&       &24:12&       &&  \\\midrule

遭逼迫时主必賜口才&       &       &13:11&21:14-15&  \\\midrule
常存忍耐就必保全靈魂&       &24:13   &13:13b&21:16-19&  \\\midrule

打仗地大震饑荒瘟疫異象神蹟&       &       &       &21:11&  \\\midrule

天國福音傳遍天下末期來到&       &24:14&13:10&&  \\\midrule
\end{tabular}%
\end{table}




















\begin{table}[H]\centering
\begin{tabular}{p{14em}|p{7.93em}|l|l|l|p{4.855em}}\toprule
 \textbf{事 件} & \textbf{地 點} & \multicolumn{1}{p{5.285em}|}{\textbf{馬太福音}} & \multicolumn{1}{p{5.57em}|}{\textbf{馬可福音}} & \multicolumn{1}{p{5.285em}|}{\textbf{路加福音}} & \textbf{約翰福音} \\  \midrule
耶路撒冷被兵圍困成荒場&       &       &       &21:20-24&  \\\midrule

那行毀壞可憎的站在...&       &24:15 &13:14a&&  \\\midrule
在猶太的/在房上的/在田裡的&       &24:16-18& 13:14b-16&&  \\\midrule
那些日懷孕的奶孩子的有禍&       &24:19&13:17 &&  \\\midrule

祈求逃走時不遇冬天安息日&       &24:20&13:18 &&  \\\midrule
空前绝后的大灾难&       &24:21&13:19 &&  \\\midrule
為選民，那日子必減少了&       &24:22&13:20  &&  \\\midrule
假基督假先知起來顯大神蹟&       &24:23-26&13:21-22&&  \\\midrule
凡事我都預先告訴要謹慎&       &&13:23&&  \\\midrule

人子降臨如閃電從東發出到西邊&       &24:27&       &&  \\\midrule

屍首在哪裡鷹也必聚在那裡&       &24:28&       &&  \\\midrule

日黑月无光星落天震動&       &24:29&13:24-25 &21:25-26a&  \\\midrule
邦國困苦人因海中波浪惊慌&       &       &       &21:26b&  \\\midrule
人子有能力榮耀駕雲降臨&       &24:30&13:26 &21:27-28&  \\\midrule

地上的萬族都要哀哭&       &24:30&       &&  \\\midrule
使者用大聲號筒從四方招聚選民&       &24:31&13:27 &&  \\\midrule

\end{tabular}%
\end{table}

\subsubsection{提醒忠仆要预备和警醒}
\begin{table}[H] \centering
\begin{tabular}{p{14em}|p{7.93em}|l|l|l|p{4.855em}} \toprule
\textbf{事 件} & \textbf{地 點} & \multicolumn{1}{p{5.285em}|}{\textbf{馬太福音}} & \multicolumn{1}{p{5.57em}|}{\textbf{馬可福音}} & \multicolumn{1}{p{5.285em}|}{\textbf{路加福音}} & \textbf{約翰福音} \\  \midrule
無花果樹和各樣樹喻神國近了&       &24:32-35&13:28-31&21:29-33&  \\\midrule
那日子時辰无人知道&       &24:36&13:32&&  \\\midrule
挪亞洪水的日子&       &24:37-39&       &&  \\\midrule
兩個人在田裡兩個女人推磨&       &24:40-41&       &&  \\\midrule

謹慎警醒常祈求见人子&       & 24:42-43 &13:33&21:34-36&  \\\midrule
你們也要預備&       & 24:44 &       &&  \\\midrule
忠心有見識的僕人&       & 24:45-47&       &&  \\\midrule

做當做的工看門的警醒不要睡著&       & &13:34-37       &&  \\\midrule

打同伴的惡僕和假冒為善人同罪&       & 24:48-51&       &&  \\\midrule

日在殿裡教訓人夜宿橄欖山&       &       &       &21:37-38&  \\\midrule

十個童女拿燈出去迎接新郎&警醒 & 25:1-13&       &&  \\\midrule
五千二千一千按才幹給銀子&有忠心 & 25:14-30&       &&  \\\midrule
人子審判萬民像分別綿羊山羊&& 25:31-46&       &&  \\\midrule
\end{tabular}%
\end{table}

\subsubsection{猶大见祭司長守殿官要銀子賣主}
\begin{table}[H] \centering
\begin{tabular}{p{14em}|p{7.93em}|l|l|l|p{4.855em}} \toprule
\textbf{事 件} & \textbf{地 點} & \multicolumn{1}{p{5.285em}|}{\textbf{馬太福音}} & \multicolumn{1}{p{5.57em}|}{\textbf{馬可福音}} & \multicolumn{1}{p{5.285em}|}{\textbf{路加福音}} & \textbf{約翰福音} \\  \midrule
\textcolor{red}{逾越節前兩天耶穌预言上十字架}&\textcolor{red}{周一}& 26:1-2&       &&  \\\midrule
祭司長守殿官商量殺耶穌&該亞法的院裡& 26:3-4&14:1-2&22:1-2&  \\\midrule
猶大见祭司長守殿官要銀子賣主& &26:14-16&14:10-11 &22:3-6&  \\\midrule
\end{tabular}%
\end{table}


\subsection{耶穌第四次預言自己受死-太26：1-2}
耶穌說完了這一切的話，就對門徒說：
你們知道，過兩天是逾越節，人子將要被交給人，釘在十字架上。















%=====================
\subsection{need more works猶大與祭司長等商議捉拿耶穌-太/可/路}
\subsubsection{馬太-21:45-46/26：3-5/26：14-16}
\paragraph{21:45-46}
祭司長和法利賽人聽見他的比喻，就看出他是指著他們說的。
他們想要捉拿他，只是怕眾人，因為眾人以他為先知.
\paragraph{26：3-5}
那時，祭司長和民間的長老聚集在大祭司稱為該亞法的院裡。大家商議要用詭計拿住耶穌，殺他，只是說：當節的日子不可，恐怕民間生亂。%（26：3-5）
\paragraph{26：14-16}
當下，十二門徒裡有一個稱為加略人猶大的，去見祭司長，說：我把他交給你們，你們願意給我多少錢？他們就給了他三十塊錢。從那時候，他就找機會要把耶穌交給他們。%（26：14-16）
\subsubsection{馬可-11:18/12:12/14：1-2/10-11}
%\subsubsection{馬可 11:18/14：1-2,10-11}
\paragraph{11:18}
祭司長和文士聽見這話，就想法子要除滅耶穌，卻又怕他，因為眾人都希奇他的教訓。
\paragraph{12:12}
他們看出這比喻是指著他們說的，就想要捉拿他，只是懼怕百姓，於是離開他走了。
\paragraph{14：1-2}
過兩天是逾越節，又是除酵節，祭司長和文士想法子怎麼用詭計捉拿耶穌，殺他。只是說：當節的日子不可，恐怕百姓生亂。
\paragraph{14：10-11}
十二門徒之中，有一個 加略人猶大去見祭司長，要把耶穌交給他們。他們聽見 就歡喜，又應許給他銀子；他就尋思如何得便把耶穌交給他們。
\subsubsection{路加-19:47b-48/22：1-6}
祭司长，和文士，与百姓的尊长，都想要杀他。 但寻不出法子来，因为百姓都侧耳听他。
除酵節，又名逾越節，近了。
祭司長和文士想法子怎麼才能殺害耶穌，是因他們懼怕百姓。
這時，撒但入了那稱為加略人猶大的心；他本是十二門徒裡的一個。
他去和祭司長並守殿官商量，怎麼可以把耶穌交給他們。
他們歡喜，就約定給他銀子。
他應允了，就找機會，要趁眾人不在跟前的時候把耶穌交給他們。
\subsubsection{综合}
\paragraph{祭司長等商議捉拿耶穌}
祭司長和文士聽見這話，就想法子要除滅耶穌，卻又怕他尋不出法子來，因為眾人都稀奇他的教訓側耳聽他。祭司長和文士看見耶穌所行的奇事，又見小孩子在殿裡喊著說「和散那歸於大衛的子孫」，就甚惱怒， 對他說：「這些人所說的，你聽見了嗎？」耶穌說：「是的。經上說『你從嬰孩和吃奶的口中完全了讚美』的話，你們沒有念過嗎？」
耶穌離開祭司長和文士，出城到伯大尼去，在那裡住宿。
\paragraph{猶大與祭司長商量賣主}
當下，十二門徒裡有一個稱為加略人猶大的，那时，撒旦进入了那称为加略人的犹大里面。去見祭司長和守卫长们商量怎样把耶稣交给他们。，說： 15 「我把他交給你們，你們願意給我多少錢？」他們聽見就歡喜，又應許給他銀子，他們就給了他三十塊錢。 16 從那時候，他就找機會要趁众人不在的时候，把耶穌交給他們。
\subsection{end OF need more works猶大與祭司長等商議捉拿耶穌-太/可/路}
%=====================

















\subsubsection{周二——逾越节晚餐}
\begin{table}[H] \centering
\begin{tabular}{p{14em}|p{7.93em}|l|l|l|p{4.855em}} \toprule
\textbf{事 件} & \textbf{地 點} & \multicolumn{1}{p{5.285em}|}{\textbf{馬太福音}} & \multicolumn{1}{p{5.57em}|}{\textbf{馬可福音}} & \multicolumn{1}{p{5.285em}|}{\textbf{路加福音}} & \textbf{約翰福音} \\  \midrule
打發彼得約翰預備逾越節筵席&\textcolor{red}{周二}耶京城里&26:17-18&14:12-15&22:7-12&  \\\midrule
使徒備逾越節筵席&耶路撒冷城里&26:19&14:16&22:13&  \\\midrule
不再吃這筵席直到在神的國裡&耶路撒冷&&&22:14-16&  \\\midrule
接過杯祝謝让大家分喝&耶路撒冷&&&22:17&  \\\midrule
不再喝這葡萄汁直等神國來到&耶路撒冷&26:29&14:25&22:18&  \\\midrule
拿起餅(身體)祝謝掰開遞門徒&耶路撒冷&26:26&14:22&22:19&  \\\midrule
飯後照樣拿起杯用血立新約&耶路撒冷&26:27-28&14:23-24&22:20&  \\\midrule

人子必死，警告賣人子的人有禍&耶路撒冷&26:20-25&14:17-18&22:21-22&  \\\midrule
門徒彼此對問是谁出卖耶稣&耶路撒冷&26:20-25&14:19-21&22:23&  \\\midrule
門徒爭論他們中間谁算為大&耶路撒冷&& &22:24&  \\\midrule

離席脫衣洗門徒腳手巾擦乾&耶路撒冷&& &&13:1-5\\\midrule
洗彼得遭拒,若不洗就與主無份&耶路撒冷&& &&13:6-9\\\midrule
你們不都是乾淨的&耶路撒冷&& &&13:10b-11\\\midrule


當彼此洗腳&耶路撒冷&& &&13:12-15\\\midrule
僕人不能大於主人&& &&&13:16-17\\\midrule
為首領的倒要像服侍人的&耶路撒冷&&&22:25-27&  \\\midrule
主许門徒坐寶座上審判12支派&耶路撒冷&& &22:28-30&  \\\midrule

再次预言有一個門徒要賣主&& &&&13:18-21\\\midrule
彼得托约翰问卖主之人&& &&&13:22-25\\\midrule
猶大吃餅立刻就出去&& &&&13:26-30\\\midrule
人子得榮耀神在人子身上得榮耀神要因自己榮耀人子&& &&&13:31-32\\\midrule
一條新命令乃是彼此相愛&& &&&13:33-35\\\midrule

\textcolor{red}{彼得}要与主同往為主捨命&& &&&13:36-37\\\midrule
主去為门徒預備地方&& &&&14:1-4\\\midrule
\textcolor{red}{多馬}问路,主是道路真理生命&& &&&14:5-7\\\midrule
\textcolor{red}{腓力}求见父,主答父在我裡面&& &&&14:8-11\\\midrule
奉主名求必成就叫父因子得榮&& &&&14:12-15\\\midrule
\end{tabular}%
% \caption{到第四个逾越节前（2）}       
 % \label{tab:addlabel}%
\end{table}


\begin{table}[H]
  \centering
    \begin{tabular}{p{14em}|p{7.93em}|l|l|l|p{4.855em}}
    \toprule
    \textbf{事 件} & \textbf{地 點} & \multicolumn{1}{p{5.285em}|}{\textbf{馬太福音}} & \multicolumn{1}{p{5.57em}|}{\textbf{馬可福音}} & \multicolumn{1}{p{5.285em}|}{\textbf{路加福音}} & \textbf{約翰福音} \\  \midrule

許賜保惠師真理聖靈&& &&&14:16-20\\\midrule
有主命又遵守的主要向他顯現&& &&&14:21\\\midrule
答\textcolor{red}{猶大}为何不向世人顯現&& &&&14:22-24\\\midrule
差來的聖靈要將一切事指教你們&& &&&14:25-26\\\midrule
主赐平安不要憂愁你们必喜樂&& &&&14:27-31\\\midrule
枝子在真葡萄樹里就多結果子&& &&&15:1-8\\\midrule
遵守主命常在主愛里彼此相愛&& &&&15:9-17\\\midrule

世人恨主者必恨信主的&& &&&15:18-25\\\midrule
保惠師來作見證门徒為主作見證&& &&&15:26-27\\\midrule
预先告诉，到时候可想起主的话&& &&&16:1-7a\\\midrule
保惠師來為罪為義為審判&& &&&16:7b-15\\\midrule
\textcolor{red}{門徒}猜往父那裡去/等不多時&& &&&16:16-18\\\midrule
门徒痛哭后憂愁要變為喜樂&& &&&16:19-22\\\midrule
向父求必得著叫你們喜樂滿足&& &&&16:23-28\\\midrule
門徒說信主是從神出來的&& &&&16:29-30\\\midrule
在世上苦難，主勝了世界&& &&&16:31-33\\\midrule
主为自己祈求&& &&&17:1-8\\\midrule
主为门徒求合一因真理成聖&& &&&17:9-19\\\midrule
主为信徒求合一榮耀愛&& &&&17:20-26\\\midrule

預言彼得三次不認主&&       &       &22:31-34&13:38\\\midrule
以后出去带錢囊和刀&       &       &       &22:35-38&  \\\midrule

唱了詩就往橄欖山去&橄欖山  &26:30 &14:26 &22:35-38&18:1 \\\midrule

告诉门徒主復活後要先去加利利&橄欖山  & &14:27-28 &22:35-38&  \\\midrule
預言彼得三次不認主&&26:31-35&  14:29-31 &      &  \\\midrule

在客西馬尼禱告&客西馬尼&26:36-46& 14:32-36 &22:39-42&  \\\midrule
天使從天上顯現加添力量&客西馬尼&       &       &22:43-44&  \\\midrule
主勉勵門徒警醒禱告&？？？ &       &14:37-41&22:45-46&  \\\midrule
猶大領兵和差役親嘴賣主&客西馬尼&26:47-50&14:42-46&22:47-48&18:2-3  \\\midrule
主說:我就是.兵丁退後倒在地上&客西馬尼&       &       &&18:4-9\\\midrule
彼得削大祭司僕人馬勒古右耳&客西馬尼&26:51-54&14:47 &22:49-51&18:10-11 \\\midrule

黑暗掌權，耶穌被捕&客西馬尼&26:55-56a&14:49&22:52-53&  \\\midrule
門徒都離開他逃走了&客西馬尼  & 26:56b &14:50-52&  \\\midrule

耶穌被捆綁到亞那院子&亞那院子&       &       &22:54&18:12-14\\\midrule
彼得一次不認主&亞那院子&26:69-70&14:66-68&22:55-57&18:15-18\\\midrule
盤問教訓时差役用手掌打主&亞那院子&&&&18:19-23\\\midrule
\end{tabular}%
% \caption{到第四个逾越节前（2）}       
 % \label{tab:addlabel}%
\end{table}


\begin{table}[H]
  \centering
    \begin{tabular}{p{14em}|p{7.93em}|l|l|l|p{4.855em}}
    \toprule
    \textbf{事 件} & \textbf{地 點} & \multicolumn{1}{p{5.285em}|}{\textbf{馬太福音}} & \multicolumn{1}{p{5.57em}|}{\textbf{馬可福音}} & \multicolumn{1}{p{5.285em}|}{\textbf{路加福音}} & \textbf{約翰福音} \\  \midrule


被转交，文士長老已經聚會&該亞法院子&26:57&14:53&&18:24\\\midrule
彼得跟进該亞法院和差役同坐&該亞法院子&26:58&14:54   &&  \\\midrule
假見證：拆毀神殿三日建造&該亞法院子&26:59-62&14:55-59 &&  \\\midrule
大祭司问主是神的兒子基督否&該亞法院子&26:63-65&14:60-64&&  \\\midrule
吐唾沫用拳頭手掌打主&該亞法院子&26:67-68&14:65&&  \\\midrule


彼得二次不認主&該亞法院子門口&26:71-72& 14:69-70a &22:58&18:25\\\midrule

過一小時彼得第三次不認主&該亞法院子 &26:73-74&14:70b-71&22:59-60&18:26-27\\\midrule

主看彼得,彼得想起主話痛哭&該亞法院子 &26:75&14:72&22:61-62&\\\midrule
耶穌被兵丁戲弄&該亞法院子&       &       &22:63-65&  \\\midrule
天亮耶穌在公會受審& \textcolor{red}{周三}公會   &27:1&15:1&22:66-71&  \\\midrule

把他捆綁交給巡撫彼拉多&彼拉多衙門& 27:2&       &&18:28-32\\\midrule

猶大悔把30塊錢給祭司長長老&殿裡&27:3-4&       &&  \\\midrule
猶大退錢遭拒出去吊死&       &27:5&       &&  \\\midrule
祭司長買窯戶的一塊田「血田」& 窯戶的田 &27:6-10&       &&  \\\midrule


彼拉多審耶稣是否猶太人的王&彼拉多衙門& 27:11-14&15:2-5&23:1-5&18:33-37\\\midrule
耶穌在希律前受審&希律府&       &       &23:6-11&  \\\midrule
希律和彼拉多从仇人成朋友&&       &       &23:12&  \\\midrule

回到彼拉多前受審&彼拉多府&27:15-18&       &23:13-17&  \\\midrule
彼拉多妻在夢中為主受苦&彼拉多府&27:19&       &&  \\\midrule
彼拉多查不出罪要釋放主&彼拉多府&&       &&18:38-39\\\midrule
眾人喊除掉耶穌釋放巴拉巴&彼拉多府&27:20-23&15:6-14 &23:18-25&18:40\\\midrule

彼拉多洗手眾人承诺血罪歸己身&彼拉多府&27:24-25&       &&  \\\midrule
釋放巴拉巴鞭打耶穌釘十字架&彼拉多府&27:26& 15:15 &&19:1 \\\midrule
巡撫全營用荊棘冠戲弄主打他&彼拉多府&27:27-31&15:16-20&&19:2-3\\\midrule

彼拉多拉戴荊棘冠的主见众人&彼拉多府外&       &       &&19:4-5\\\midrule
祭司長差役喊釘他十字架&彼拉多府外&       &       &&19:6-7\\\midrule
彼拉多惊闻主是神子继续审主&彼拉多府&       &       &&19:8-11\\\midrule
彼拉多要放主被猶太人威胁&彼拉多府外&       &       &&19:12\\\midrule
犹太人继续迫胁彼拉多钉主&鋪華石處&       &       &&19:13-15\\\midrule
將耶穌交給他們去釘十字架&鋪華石處&       &       &&19:16\\\midrule

古利奈人西門背十架跟隨耶穌&       &27:32&15:21-22 &23:26&  \\\midrule
耶穌劝婦女們为自己兒女哀哭&       &       &       &23:27-31&  \\\midrule

兵丁拿苦膽沒藥調和酒給主喝&髑髏地&27:33-34&15:23 &&  \\\midrule

巳初兩犯人与主同釘十架&髑髏地&27:38&15:25,27&23:32-33&19:17-18\\\midrule
\textcolor{red}{耶稣求父赦免釘他之人}&       &   &      &23:34a&  \\\midrule
釘完十字架，兵丁拈鬮分衣服&髑髏地&27:35-36&15:24&23:34b&19:23-34\\\midrule

\end{tabular}%
% \caption{到第四个逾越节前（2）}       
 % \label{tab:addlabel}%
\end{table}


\begin{table}[H]
  \centering
    \begin{tabular}{p{14em}|p{7.93em}|l|l|l|p{4.855em}}
    \toprule
    \textbf{事 件} & \textbf{地 點} & \multicolumn{1}{p{5.285em}|}{\textbf{馬太福音}} & \multicolumn{1}{p{5.57em}|}{\textbf{馬可福音}} & \multicolumn{1}{p{5.285em}|}{\textbf{路加福音}} & \textbf{約翰福音} \\  \midrule

主吩咐约翰"\textcolor{red}{看你的母親}"&       & &        &&19:25-27\\\midrule

百姓站在那裡觀看&       & &        &23:35a&  \\\midrule
經過的人譏誚神兒子下十字架&       &27:39-40&15:29-30&&  \\\midrule
祭司長文士長老戲弄他&       &27:41-43&15:31-32a&&  \\\midrule


官府也嗤笑&       & &        &23:35b&  \\\midrule

牌子寫著“這是猶太人的王”& &27:37&15:26&23:38&19:19-22\\\midrule
一犯人譏誚主& &27:44&15:32b &23:39&  \\\midrule
一个犯人求主記念& &&       &23:40-42&  \\\midrule
\textcolor{red}{主应许犯人今日同主在樂園}& &&       &23:43&  \\\midrule

午正到申初日黑了（幔裂）&       &27:45& 15:33 &23:44-45&  \\\midrule
\textcolor{red}{申初耶穌喊我的神为何離棄我}&       &27:46&15:34&&  \\\midrule

為要使經上的話應驗就說\textcolor{red}{我渴了}&       & &        &&19:28-29\\\midrule
兵丁拿醋戲弄主&  &27:48&15:36a  &23:36-37&  \\\midrule

站在那的人以为呼叫以利亞&       &27:47& 15:35  &&  \\\midrule
其餘的人看以利亞來救他不來&       &27:49&15:36b  &&  \\\midrule
耶穌嘗了那醋就說\textcolor{red}{成了}&       &&       &&19:30\\\midrule
\textcolor{red}{耶穌將靈魂交父手}&       &&       &23:46a&  \\\midrule
耶穌氣就斷了&       &27:50&15:37 &23:46b&  \\\midrule
地震石崩墓開已睡聖徒復活&       &27:51-52&       &&  \\\midrule
殿裡幔子從上到下裂為兩半&       && 15:38      &&  \\\midrule
百夫長歸榮耀於神&       &27:54 &15:39 &23:47&  \\\midrule
眾人捶著胸回去&       &       &       &23:48&  \\\midrule
耶穌熟人和婦女站著看這些事&       &27:55-56&15:40-41 &23:49&  \\\midrule

猶太人求彼拉多打斷犯人的腿&預備日\textcolor{red}{周三} && &&19:31-33\\\midrule
一個兵拿槍扎主肋有血水流出&預備日\textcolor{red}{周三}&& &&19:34-37\\\midrule

約瑟求領耶穌的身體&預備日 && &&19:38\\\midrule
尼哥迪慕帶沒藥沉香一百斤&預備日\textcolor{red}{周三}&& &&19:39\\\midrule


約瑟葬耶穌婦女在旁观看&\textcolor{red}{周三}&27:57-61 &  15:42-47&23:50-54&19:40-42\\\midrule

祭司長法利賽人求把守坟墓\textcolor{red}{周四}&預備日第二天&27:62-66&&&\multicolumn{1}{p{6.85em}}{\textbf{19:31大安息日}}  \\\midrule
婦女回去預備香料香膏\textcolor{red}{周五}&过安息日第二天     &&       &23:55&  \\\midrule

她们在安息日便遵着诫命安息&\textcolor{red}{周六}&&       &23:56&  \\\midrule
天快亮两个馬利亞來看墳墓&\textcolor{red}{周日清晨}&28:1&       &&20:1 \\\midrule
地大震，见到主使者&       &28:2-3&       &&  \\\midrule
看守的人嚇得渾身亂戰&       &28:4&       &&  \\\midrule
使者安慰妇女，并告知主复活&       &28:5-7&       &&  \\\midrule
兵丁受銀錢造谣说门徒偷尸&       &28:11-15 &       &&  \\\midrule



\end{tabular}%
% \caption{到第四个逾越节前（2）}       
 % \label{tab:addlabel}%
\end{table}


\begin{table}[H]
  \centering
    \begin{tabular}{p{14em}|p{7.93em}|l|l|l|p{4.855em}}
    \toprule
    \textbf{事 件} & \textbf{地 點} & \multicolumn{1}{p{5.285em}|}{\textbf{馬太福音}} & \multicolumn{1}{p{5.57em}|}{\textbf{馬可福音}} & \multicolumn{1}{p{5.285em}|}{\textbf{路加福音}} & \textbf{約翰福音} \\  \midrule
抹大拉馬利亞告訴使徒墳墓空&       &28:8&       &&20:2 \\\midrule

彼得约翰去墳墓&       &&       &&20:3-5 \\\midrule


彼得見空墳心裡稀奇&       &       &       &24:12&20:6-7\\\midrule
约翰进去看墳墓信墳空&       &&       &&20:8\\\midrule
彼得约翰不懂死裡復活回去了&       &&       &&20:9-10\\\midrule

抹大拉馬利亞哭泣遇天使&       &&       &&20:11-13\\\midrule
抹大拉馬利亞遇見耶穌&       && 16:9&&20:14-17\\\midrule

婦女們黎明時見空墳&       &       &16:1-4&24:1-7&  \\\midrule
婦女們見白袍少年人&       &       &16:5-8&&  \\\midrule
婦女們遇見耶穌&       &28:9-10&       &24:8-11&  \\\midrule
婦女們告訴使徒耶穌復活&       &16:10-11&       &24:8-11&20:18 \\\midrule

主果然復活現給西門&以馬忤斯路上&       &       &24:34&  \\\midrule
已睡聖徒復活進聖城向人顯現&& 27:53&       &&  \\\midrule

在以馬忤斯路上向兩門徒顯現&以馬忤斯路上&       &16:12-13 &24:13-35&  \\\midrule
耶穌向其他門徒顯現&耶路撒冷 &       &16:14-18&24:36-49&20:19-21\\\midrule

主向他們吹氣受聖靈赦罪&耶路撒冷 &       &&&20:22-23 \\\midrule
多馬不信主已经复活&耶路撒冷 &       &&&20:24-25\\\midrule

過八日主向多馬和門徒显现&耶路撒冷 &       &&&20:26-29\\\midrule
許多神蹟沒有記在這書上& &       &&&20:30-31\\\midrule

彼得等六人一夜沒打著打鱼&提比哩亞海邊 &       &&&21:1-3\\\midrule
主在岸边指示撒網于船右&提比哩亞海邊 &       &&&21:4-8\\\midrule

岸上有炭火有魚又有餅&提比哩亞海邊 &       &&&21:9-13\\\midrule
主三次邀彼得爱主喂养主羊&提比哩亞海邊 &       &&&21:14-17\\\midrule

呼召彼得跟从主&提比哩亞海邊 &       &&&21:18-23\\\midrule
此书的目为主作見證&&       &&&21:24-25\\\midrule

十一個門徒往加利利领受使命&約定的山&28:16-20 &       &&  \\\bottomrule
耶穌升天&伯大尼 &       & 16:19-20&24:50-53&  \\\midrule
四十天之顯現給使徒&徒1:3-11 & &&&\\\midrule
一時顯給五百多弟兄&林前15:3-8 & &&&\\\bottomrule

    \end{tabular}%
% \caption{到第四个逾越节前（2）}       
  \label{tab:addlabel}%
\end{table}%






